\section*{Chapitre \ref{ch:m1:index-and-single-stock-futures-worldwide} --- Contrats à terme sur indices actions et sur actions dans le monde}

\begin{solution}{exo:m1:index-and-single-stock-futures-worldwide:1}
300\,000~\$, 54\,000~\euro et HK\$1\,300\,000~; un mouvement de 1~\% vaut 3\,000~\$, 540~\euro et
HK\$13\,000 par contrat.
\end{solution}

\begin{solution}{exo:m1:index-and-single-stock-futures-worldwide:2}
$5\,400 \times (1 + 0.025 \times 0.25) - 5\,412.0 = 5\,433.75 - 5\,412.0 = 21.75$
points.
\end{solution}

\begin{solution}{exo:m1:index-and-single-stock-futures-worldwide:3}
$(158.2 - 171.5) \times \text{\euro}100 \times 50 = -\text{\euro}66\,500$.
\end{solution}

\begin{solution}{exo:m1:index-and-single-stock-futures-worldwide:4}
$+2.08\%$, $-1.46\%$, $-1.78\%$, $-0.90\%$, $-0.61\%$~; sur cinq ans
$(163.5/168.0)^{1/5} - 1 = -0.54\%$ par an. Le strip est un prix, non une
prévision~: c'est ce que les vendeurs de risque de dividende, surtout des hedgers de
produits structurés, acceptent de s'en défaire. L'écart avec la
prévision de l'analyste est la prime que touche l'acheteur pour porter ce risque, ainsi que la
possibilité que l'analyste se trompe.
\end{solution}

\begin{solution}{exo:m1:index-and-single-stock-futures-worldwide:5}
Les trois séances de Hong Kong se situent dans les plages horaires des deux autres~: $2.75 + 3.5
+ 10 = 16.25$ heures avec les trois ouvertes. Rien n'est ouvert de 21:00 à 22:00
UTC, la pause quotidienne du contrat américain. En décembre, Francfort et Chicago
sont une heure plus tard en UTC (leurs horloges reculent~; l'UTC et Hong Kong ne
bougent pas)~: la plage européenne va de 01:10 à 21:00 et la pause tombe à
22:00.
\end{solution}

\begin{solution}{exo:m1:index-and-single-stock-futures-worldwide:6}
$200 \times \text{\euro}10 \times 5\,400 \times 0.0015 \times 2 =
+\text{\euro}32\,400$. La position longue paie le spread comme financement~; ayant figé 40
lorsque le marché facture désormais 55, la position vaut la différence pour
les deux années restantes.
\end{solution}

\begin{solution}{exo:m1:index-and-single-stock-futures-worldwide:7}
Le 16 décembre, l'heure vide va de 22:00 à 23:00 UTC, contre 21:00 à 22:00
en septembre. Pendant la semaine du 14 septembre, l'ensemble change à chaque
frontière de séance~: dimanche 22:00 UTC (ouverture du contrat américain), puis chaque
jour 00:10 (ouverture de l'Europe), 01:15, 04:00, 05:00, 08:30, 09:00 et 19:00 (Hong
Kong, trois séances), 20:00 (fermeture de l'Europe), 21:00 et 22:00 (la
pause américaine), jusqu'au vendredi 21:00, quand le dernier contrat clôture.
\end{solution}

\begin{solution}{exo:m1:index-and-single-stock-futures-worldwide:8}
La statistique compte les contrats, et ceux-ci diffèrent en taille par des facteurs de
mille. Le total 2025 recule de 42~\% alors que le volume des contrats à terme progresse de 8.6~\%~: le décompte
est dominé par des options sur actions de très petite taille dans quelques marchés, et évolue
avec les changements réglementaires touchant ces contrats, et non avec le montant de risque
transféré. Une société devrait classer les marchés par notionnel ou, mieux encore, par la
revenus accessibles à son type de stratégie~: volume en valeur, spreads, frais,
accès et concurrence. Certains des marchés derrière les 63~\% sont aussi fermés,
ou presque, aux sociétés de trading étrangères pour compte propre.
\end{solution}

\begin{solution}{pb:m1:index-and-single-stock-futures-worldwide:1}
\textbf{1.} $171.5/5\,400 = 3.18\%$.
\textbf{2.} $+2.08\%$.
\textbf{3.} $-0.54\%$ par an.
\textbf{4.} $168.0 \times 1.04^5 = 204.4$~; le contrat à terme à 163.5 est inférieur de 20~\%.
\textbf{5.} Les dealers qui couvrent l'exposition longue aux dividendes que les produits structurés
leur laissent~; ils vendent à décote, car personne d'autre ne peut servir de
contrepartie, en taille, à un vendeur naturel pour une échéance à cinq ans.
\textbf{6.} $200 \times 163.5 \times 100 = \text{\euro}3.27$ millions d'euros.
\textbf{7.} $(204.4 - 163.5) \times 20\,000 = +\text{\euro}818\,000$.
\textbf{8.} $168.0 \times 0.7 = 117.6$~; $(117.6 - 163.5) \times 20\,000 =
-\text{\euro}918\,000$.
\textbf{9.} $(110 - 163.5) \times 20\,000 = -\text{\euro}1.07$ million d'euros, en
espèces, dans le mois, sur une position de 3.27 millions d'euros. Avoir raison en 2031
suppose de survivre aux \omterm{def:m1:financing:margincall}{appels de marge} traités au \cref{ch:m1:margin} jusqu'alors.
\textbf{10.} $20\% \times 171.5 = 34.3$ points.
\textbf{11.} La fair value du contrat à terme augmente de 34.3 points~: $500 \times 10
\times 34.3 = +\text{\euro}171\,500$. Une position longue sur contrat à terme est courte en dividendes, donc une
baisse de dividende est un gain.
\textbf{12.} Acheter 50 \omterm{def:m1:index-and-single-stock-futures-worldwide:divfut}{contrats à terme sur dividendes} de 2027~: $50 \times 100 \times 34.3 =
\text{\euro}171\,500$.
\textbf{13.} Les produits qu'elle vend versent aux investisseurs le rendement de prix de l'indice,
souvent avec une garantie~; le desk se couvre en détenant l'indice, par le biais d'actions
ou des contrats à terme, et encaisse les dividendes sans en devoir aucun. Son profit dépend
des dividendes pendant la durée de vie du produit, de cinq à dix ans.
\textbf{14.} Le \omterm{def:m1:index-and-single-stock-futures-worldwide:trf}{contrat à terme de rendement total} verse les dividendes réalisés au lieu de
intégrer une estimation de ceux-ci, si bien que la couverture n'a aucune exposition aux
dividendes au départ, et fixe le \omterm{def:m1:delta-one-instruments:swap}{spread de financement} pour la durée de vie du produit.
\textbf{15.} Non. C'est le prix auquel le risque de dividende se fixe entre
vendeurs structurels et les capitaux prêts à le porter, et il varie avec l'appétit
de ces capitaux autant qu'avec les anticipations de dividendes.
\textbf{16.} Le fait que les détenteurs soient des investisseurs à effet de levier ou évalués en valeur de marché, qui
sont contraints de sortir lorsque les cours baissent et que les marges augmentent, si bien que le contrat
se comporte en crise comme une position à effet de levier sur actions, et non comme la
variable lente sur laquelle il se règle.
\textbf{17.} Rien de négociable~: les dividendes de 2031 n'existent pas encore. Le
contrat ne peut être couvert qu'approximativement, avec l'indice, avec d'autres
points du strip, ou avec des options~; son risque est en grande partie non couvrable, ce qui
explique pourquoi il porte une prime.
\textbf{18.} Moins liquides. Ils existent pour la même raison au niveau d'une
action~: les desks d'options et de produits structurés détiennent des dividendes sur actions individuelles
et ont besoin de les vendre, et une baisse de dividende sur un seul titre est un
risque important et discontinu.
\textbf{19.} \textbf{$+2.1\%$ pour la première année~; $-0.5\%$ par an sur cinq
ans.}
\textbf{20.} La composante dividende de la détention d'actions, à partir de son cours et de son
financement.
\end{solution}

\begin{solution}{iq:m1:index-and-single-stock-futures-worldwide:1}
Le E-mini S\&P 500 à Chicago, le contrat à terme sur l'EURO STOXX 50 à Francfort, les
contrats Hang Seng ou Nikkei en Asie. Le E-mini expire le troisième vendredi
contre une cotation d'ouverture~; le contrat sur l'EURO STOXX 50 le même jour à midi
contre une moyenne sur dix minutes~; le contrat Hang Seng est mensuel,
et expire l'avant-dernier jour de séance, contre une moyenne de cotations sur cinq minutes.
\iqlookfor{un détail concret sur l'échéance de chaque contrat.}
\end{solution}

\begin{solution}{iq:m1:index-and-single-stock-futures-worldwide:2}
Un contrat à terme qui se règle sur les dividendes qu'un indice verse sur
une période, en \omterm{def:m1:index-and-single-stock-futures-worldwide:point}{points d'indice}. Vendeurs~: les dealers en position longue de dividendes
issus de produits structurés et de books d'options. Acheteurs~: les fonds qui captent la
décote, et les traders de valeur relative face au contrat sur indice, aux options et aux dividendes sur titres individuels.
\iqlookfor{le vendeur structurel.}
\end{solution}

\begin{solution}{iq:m1:index-and-single-stock-futures-worldwide:3}
Si le dividende baisse avant l'échéance de mon contrat à terme, je gagne~: le contrat était
coté sous le spot majoré des intérêts, de la valeur du dividende attendu~; l'indice ne
baissera plus de ce montant à la \omterm{def:m1:shares-corporate-actions-indices:dividend}{date de détachement}, donc la fair value monte du
dividende, exprimé en \omterm{def:m1:index-and-single-stock-futures-worldwide:point}{points d'indice}. L'indice au comptant peut lui-même baisser à l'annonce~;
il s'agit d'un mouvement de prix distinct et ordinaire.
\iqlookfor{position longue sur contrat à terme équivaut à position courte sur dividendes, distincte de la
réaction du cours.}
\end{solution}

\begin{solution}{iq:m1:index-and-single-stock-futures-worldwide:4}
Le \omterm{def:m1:delta-one-instruments:swap}{spread de financement} à terme sur actions, seul. Le contrat à terme ordinaire regroupe une
estimation des dividendes et un taux de financement pour une échéance donnée~; le \omterm{def:m1:index-and-single-stock-futures-worldwide:trf}{contrat à terme
de rendement total} verse les dividendes réalisés et facture un taux overnight réalisé, est
coté en spread et liste des échéances jusqu'à dix ans, de sorte qu'une vue ou une couverture
sur le repo d'actions puisse être tenue comme position compensée et margée.
\iqlookfor{«~coté en spread~» et le motif de bilan bancaire.}
\end{solution}

\begin{solution}{iq:m1:index-and-single-stock-futures-worldwide:5}
Séances à l'heure locale de la bourse, avec un nom de fuseau IANA, jamais des décalages UTC~;
plusieurs séances par produit et par jour, pouvant franchir minuit~; les jours de semaine
où chaque séance ouvre~; une table des jours fériés et des demi-journées par bourse et par année~;
les exceptions du jour d'échéance tirées du référentiel des contrats~; une plage de dates d'effet
sur chaque ligne, car les horaires changent. Toutes les requêtes sont traitées en UTC
par une seule fonction que tous les systèmes partagent.
\iqlookfor{noms de fuseau et non décalages~; dates d'effet.}
\end{solution}

\begin{solution}{iq:m1:index-and-single-stock-futures-worldwide:6}
Convertir les contrats en notionnel avec le multiplicateur et le prix de chaque contrat, et
pour les options, en prime ou en notionnel ajusté du delta. Séparer les contrats à terme des
options. Ajouter ensuite ce qui détermine si une société peut y gagner quoi que ce soit~:
le spread typique en points de base, la contrainte de tick, les frais et la taxe par transaction, qui
est autorisé à négocier et par quel intermédiaire, et la part des volumes
professionnels face aux particuliers. Le classement qui en résulte ressemble peu
au palmarès habituel.
\iqlookfor{conversion en notionnel d'abord, puis économie et accès.}
\end{solution}
